% UC: Lista Opere
\begin{table}[H]
    \begin{tabular}{|p{\linewidth}|}
        \hline
        \cellcolor{gray!100}
        \color{white}
        \centerline{\textbf{Caso d'uso:} Lista Opere} \\
        \hline
        \textbf{ID:} 1 \\
        \hline
        \cellcolor{gray!20}
        \textbf{Breve descrizione:} \\
        \cellcolor{gray!20}
        L'utente ottiene la lista di opere memorizzate nel sistema. \\
        \hline
        \textbf{Attori primari:} \\
        \begin{minipage}{\linewidth}
            Cliente \\
            Biglietteria \\
            Amministratore
        \end{minipage}
        \vspace{0pt} \\  % Utilizzo \vspace{} per regolare lo spazio tra le sezioni delle tabelle.
        \hline
        \textbf{Attori secondari:} \\
        Nessuno. \\
        \hline
        \cellcolor{gray!20}
        \textbf{Precondizioni:} \\
        \cellcolor{gray!20}
        \begin{minipage}{\linewidth}
            \begin{enumerate}
                \item L'utente è posizionato nella sezione informazioni.
            \end{enumerate}
        \end{minipage} \\
        \hline
        \textbf{Sequenza degli eventi principale:}
        \begin{enumerate}
            \item Il caso d'uso inizia quando l'utente entra nella sezione informazioni.
            \item \textbf{If} nel sistema è memorizzata almeno un'opera
            \begin{enumerate}
                \item \textbf{For} ogni opera nel sistema
                \begin{enumerate}
                    \item[] \textbf{punto di estensione:} filtro ricerca
                    \item Vengono visualizzate le informazioni essenziali dell'opera.
                \end{enumerate}
            \end{enumerate}
            \item \textbf{Else}
            \begin{enumerate}
                \item Il sistema informa l'utente che non vi sono opere memorizzate.
            \end{enumerate}
        \end{enumerate} \\
        \hline
        \cellcolor{gray!20}
        \textbf{Postcondizioni:} \\
        \cellcolor{gray!20}
        \begin{minipage}{\linewidth}
            \begin{enumerate}
                \item Viene visualizzata a schermo la lista delle opere memorizzate.
            \end{enumerate}
        \end{minipage} \\
        \hline
        \textbf{Sequenza degli eventi alternativa:} \\
        Nessuna. \\
        \hline
    \end{tabular}
\end{table}

% UC EXT: Ricerca Opera
\begin{table}[H]
    \begin{tabular}{|p{\linewidth}|}
        \hline
        \cellcolor{gray!100}
        \color{white}
        \centerline{\textbf{Caso d'uso di estensione:} Ricerca Opera} \\
        \hline
        \textbf{ID:} 2 \\
        \hline
        \cellcolor{gray!20}
        \textbf{Breve descrizione:} \\
        \cellcolor{gray!20}
        \textbf{Segmento 1:} L'utente filtra la lista di opere tramite parola chiave ricercata nel nome. \\
        \hline
        \textbf{Attori primari:} \\
        \begin{minipage}{\linewidth}
            Cliente \\
            Biglietteria \\
            Amministratore
        \end{minipage}
        \vspace{0pt} \\  % Utilizzo \vspace{} per regolare lo spazio tra le sezioni delle tabelle.
        \hline
        \textbf{Attori secondari:} \\
        Nessuno. \\
        \hline
        \cellcolor{gray!20}
        \textbf{Precondizioni del segmento 1:} \\
        \cellcolor{gray!20}
        \begin{minipage}{\linewidth}
            \begin{enumerate}
                \item L'utente ha inserito una chiave di ricerca.
                \item L'utente ha premuto il pulsante di ricerca.
            \end{enumerate}
        \end{minipage} \\
        \hline
        \textbf{Sequenza degli eventi principale del segmento 1:}
        \begin{enumerate}
            \item Le opere che non contengono la parola chiave vengono eliminate dalla lista risultante.
        \end{enumerate} \\
        \hline
        \cellcolor{gray!20}
        \textbf{Postcondizioni del segmento 1:} \\
        \cellcolor{gray!20}
        \begin{minipage}{\linewidth}
            \begin{enumerate}
                \item La lista risultante è filtrata.
            \end{enumerate}
        \end{minipage} \\
        \hline
    \end{tabular}
\end{table}

% UC: Visualizza Opera
\begin{table}[H]
    \begin{tabular}{|p{\linewidth}|}
        \hline
        \cellcolor{gray!100}
        \color{white}
        \centerline{\textbf{Caso d'uso:} Visualizza Opera} \\
        \hline
        \textbf{ID:} 3 \\
        \hline
        \cellcolor{gray!20}
        \textbf{Breve descrizione:} \\
        \cellcolor{gray!20}
        L'utente ottiene le informazioni anagrafiche riguardo una singola opera. \\
        \hline
        \textbf{Attori primari:} \\
        \begin{minipage}{\linewidth}
            Cliente \\
            Biglietteria \\
            Amministratore
        \end{minipage}
        \vspace{0pt} \\ % Utilizzo \vspace{} per regolare lo spazio tra le sezioni delle tabelle.
        \hline
        \textbf{Attori secondari:} \\
        Nessuno. \\
        \hline
        \cellcolor{gray!20}
        \textbf{Precondizioni:} \\
        \cellcolor{gray!20}
        \begin{minipage}{\linewidth}
            \begin{enumerate}
                \item L'utente è posizionato nella sezione informazioni.
            \end{enumerate}
        \end{minipage} \\
        \hline
        \textbf{Sequenza degli eventi principale:}
        \begin{enumerate}
            \item Il caso d'uso inizia quando l'utente richiede le informazioni di un'opera.
            \item Il sistema cerca le informazioni anagrafiche nella memoria.
            \item Il sistema mostra le informazioni a schermo.
            \item \textbf{include}(Lista Regie)
            \item[] \textbf{punto di estensione:} modifica
        \end{enumerate} \\
        \hline
        \cellcolor{gray!20}
        \textbf{Postcondizioni:} \\
        \cellcolor{gray!20}
        \begin{minipage}{\linewidth}
            \begin{enumerate}
                \item Sono visualizzate a schermo le informazioni anagrafiche dell'opera.
            \end{enumerate}
        \end{minipage} \\
        \hline
        \textbf{Sequenza degli eventi alternativa:} \\
        Nessuna. \\
        \hline
    \end{tabular}
\end{table}

% UC: Lista Regie
\begin{table}[H]
    \begin{tabular}{|p{\linewidth}|}
        \hline
        \cellcolor{gray!100}
        \color{white}
        \centerline{\textbf{Caso d'uso:} Lista Regie} \\
        \hline
        \textbf{ID:} 4 \\
        \hline
        \cellcolor{gray!20}
        \textbf{Breve descrizione:} \\
        \cellcolor{gray!20}
        L'utente ottiene la lista di regie relative all'opera che sta visualizzando. \\
        \hline
        \textbf{Attori primari:} \\
        \begin{minipage}{\linewidth}
            Cliente \\
            Biglietteria \\
            Amministratore
        \end{minipage}
        \vspace{0pt} \\ % Utilizzo \vspace{} per regolare lo spazio tra le sezioni delle tabelle.
        \hline
        \textbf{Attori secondari:} \\
        Nessuno. \\
        \hline
        \cellcolor{gray!20}
        \textbf{Precondizioni:} \\
        \cellcolor{gray!20}
        \begin{minipage}{\linewidth}
            \begin{enumerate}
                \item L'utente sta visualizzando una singola opera.
            \end{enumerate}
        \end{minipage} \\
        \hline
        \textbf{Sequenza degli eventi principale:}
        \begin{enumerate}
            \item Il caso d'uso inizia quando l'utente visualizza una singola opera.
            \item \textbf{If} per l'opera è memorizzata almeno una regia
            \begin{enumerate}
                \item \textbf{For} ogni regia dell'opera
                \begin{enumerate}
                    \item Il sistema mostra le informazioni della regia a schermo.
                \end{enumerate}
            \end{enumerate}
            \item \textbf{Else}
            \begin{enumerate}
                \item Il sistema informa l'utente che non vi sono regie memorizzate.
            \end{enumerate}
        \end{enumerate} \\
        \hline
        \cellcolor{gray!20}
        \textbf{Postcondizioni:} \\
        \cellcolor{gray!20}
        \begin{minipage}{\linewidth}
            \begin{enumerate}
                \item E' visualizzata a schermo la lista delle regie associate all'opera.
            \end{enumerate}
        \end{minipage} \\
        \hline
        \textbf{Sequenza degli eventi alternativa:} \\
        Nessuna. \\
        \hline
    \end{tabular}
\end{table}

% UC: Lista Spettacoli
\begin{table}[H]
    \begin{tabular}{|p{\linewidth}|}
        \hline
        \cellcolor{gray!100}
        \color{white}
        \centerline{\textbf{Caso d'uso:} Lista Spettacoli} \\
        \hline
        \textbf{ID:} 5 \\
        \hline
        \cellcolor{gray!20}
        \textbf{Breve descrizione:} \\
        \cellcolor{gray!20}
        L'utente visualizzare l'elenco degli spettacoli disponibili presso il teatro. \\
        \hline
        \textbf{Attori primari:} \\
        \begin{minipage}{\linewidth}
            Cliente \\
            Biglietteria \\
            Amministratore \\
        \end{minipage}
        \vspace{-10pt} \\ 
        \hline
        \textbf{Attori secondari:} \\
        Nessuno. \\
        \hline
        \cellcolor{gray!20}
        \textbf{Precondizioni:} \\
        \cellcolor{gray!20}
        \begin{minipage}{\linewidth}
            \begin{enumerate}
                \item L'utente è posizionato nella sezione informazione.
            \end{enumerate}
        \end{minipage} \\
        \hline
        \textbf{Sequenza degli eventi principale:}
        \begin{enumerate}
            \item Il caso d'uso inizia quando l'utente richiede di visualizzare la lista completa degli spettacoli.
            \item \textbf{If} c'è almeno uno spettacolo memorizzato nel Sistema
            \begin{enumerate}
                \item \textbf{For} ogni spettacolo nel Sistema
                \begin{enumerate}
                    \item Il sistema mostra un'anteprima dello spettacolo con informazioni basiche (es. titolo, data, genere). % Disponibilità?
                \end{enumerate}
                \item[] \textbf{punto di estensione:} filtrareElenco
            \end{enumerate}
            \item \textbf{Else}
            \begin{enumerate}
                \item Il Sistema informa l'utente che non vi sono spettacoli memorizzati.
            \end{enumerate}
        \end{enumerate} \\
        \hline
        \cellcolor{gray!20}
        \textbf{Postcondizioni:} \\
        \cellcolor{gray!20}
        \begin{minipage}{\linewidth}
            \begin{enumerate}
                \item Viene visualizzata a schermo la lista degli spettacoli memorizzati.
            \end{enumerate}
        \end{minipage} \\
        \hline
        \textbf{Sequenza degli eventi alternativa:} \\
        Nessuna. \\
        \hline
    \end{tabular}
\end{table}

% UC EXT: Ricerca Spettacolo
\begin{table}[H]
    \begin{tabular}{|p{\linewidth}|}
        \hline
        \cellcolor{gray!100}
        \color{white}
        \centerline{\textbf{Caso d'uso di estensione:} Ricerca Spettacolo} \\
        \hline
        \textbf{ID:} 6 \\
        \hline
        \cellcolor{gray!20}
        \textbf{Breve descrizione:} \\
        \cellcolor{gray!20}
        \textbf{Segmento 1:} L'utente filtra la lista di spettacoli applicando eventuali filtri (es. titolo, data, genere). \\
        \hline
        \textbf{Attori primari:} \\
        \begin{minipage}{\linewidth}
            Cliente \\
            Biglietteria \\
            Amministratore \\
        \end{minipage}
        \vspace{-10pt} \\ 
        \hline
        \textbf{Attori secondari:} \\
        Nessuno. \\
        \hline
        \cellcolor{gray!20}
        \textbf{Precondizioni del segmento 1:} \\
        \cellcolor{gray!20}
        \begin{minipage}{\linewidth}
            \begin{enumerate}
                \item L'utente ha inserito i criteri di ricerca.
                \item L'utente ha premuto il pulsante di ricerca.
            \end{enumerate}
        \end{minipage} \\
        \hline
        \textbf{Sequenza degli eventi principale del segmento 1:}
        \begin{enumerate}
            \item Il Sistema ricerca, nell'elenco di spettacoli, degli spettacoli coincidenti.
            \item \textbf{If} il Sistema trova almeno uno spettacolo coincidente
            \begin{enumerate}
                \item \textbf{For} ogni spettacolo coincidente
                \begin{enumerate}
                    \item Il sistema mostra un'anteprima dello spettacolo con informazioni basiche.
                \end{enumerate}
            \end{enumerate}
            \item \textbf{Else}
            \begin{enumerate}
                \item Il Sistema informa l'utente che non vi sono spettacoli coincidenti.
            \end{enumerate}
        \end{enumerate} \\
        \hline
        \cellcolor{gray!20}
        \textbf{Postcondizioni del segmento 1:} \\
        \cellcolor{gray!20}
        \begin{minipage}{\linewidth}
            \begin{enumerate}
                \item L'utente visualizza l'elenco filtrato.
            \end{enumerate}
        \end{minipage} \\
        \hline
    \end{tabular}
\end{table}

% UC: Visualizza Spettacolo
\begin{table}[H]
    \begin{tabular}{|p{\linewidth}|}
        \hline
        \cellcolor{gray!100}
        \color{white}
        \centerline{\textbf{Caso d'uso:} Visualizza Spettacolo} \\
        \hline
        \textbf{ID:} 7 \\
        \hline
        \cellcolor{gray!20}
        \textbf{Breve descrizione:} \\
        \cellcolor{gray!20}
        L'utente visualizza i dettagli di uno spettacolo selezionato dalla lista. \\
        \hline
        \textbf{Attori primari:} \\
        \begin{minipage}{\linewidth}
            Cliente \\
            Biglietteria \\
            Amministratore \\
        \end{minipage}
        \vspace{-10pt} \\ 
        \hline
        \textbf{Attori secondari:} \\
        Nessuno. \\
        \hline
        \cellcolor{gray!20}
        \textbf{Precondizioni:} \\
        \cellcolor{gray!20}
        \begin{minipage}{\linewidth}
            \begin{enumerate}
                \item L'utente visualizza la lista degli spettacoli.
                \item L'utente ha selezionato uno spettacolo della lista.
            \end{enumerate}
        \end{minipage} \\
        \hline
        \textbf{Sequenza degli eventi principale:}
        \begin{enumerate}
            \item Il Sistema mostra tutti i dettagli dello spettacolo: titolo, descrizione, durata, ecc.
            \item \textbf{include}(Lista Eventi)
            \item \textbf{If} l'utente è autenticato come Amministratore
            \begin{enumerate}
                \item[] \textbf{punto di estensione:} modificaSpettacolo
            \end{enumerate}
        \end{enumerate} \\
        \hline
        \cellcolor{gray!20}
        \textbf{Postcondizioni:} \\
        \cellcolor{gray!20}
        \begin{minipage}{\linewidth}
            \begin{enumerate}
                \item Sono visualizzate a schermo le informazioni anagrafiche dello spettacolo.
            \end{enumerate}
        \end{minipage} \\
        \hline
        \textbf{Sequenza degli eventi alternativa:} \\
        Nessuna. \\
        \hline
    \end{tabular}
\end{table}

% UC: Lista Eventi
\begin{table}[H]
    \begin{tabular}{|p{\linewidth}|}
        \hline
        \cellcolor{gray!100}
        \color{white}
        \centerline{\textbf{Caso d'uso:} Lista Eventi} \\
        \hline
        \textbf{ID:} 8 \\
        \hline
        \cellcolor{gray!20}
        \textbf{Breve descrizione:} \\
        \cellcolor{gray!20}
        L'utente visualizza l'elenco di tutti gli eventi programmati per uno spettacolo. \\
        \hline
        \textbf{Attori primari:} \\
        \begin{minipage}{\linewidth}
            Cliente \\
            Biglietteria \\
            Amministratore \\
        \end{minipage}
        \vspace{-10pt} \\ 
        \hline
        \textbf{Attori secondari:} \\
        Nessuno. \\
        \hline
        \cellcolor{gray!20}
        \textbf{Precondizioni:} \\
        \cellcolor{gray!20}
        \begin{minipage}{\linewidth}
            \begin{enumerate}
                \item L'utente visualizza i dettagli di uno spettacolo.
            \end{enumerate}
        \end{minipage} \\
        \hline
        \textbf{Sequenza degli eventi principale:}
        \begin{enumerate}
            \item Il caso d'uso inizia dopo che l'utente visualizza i dettagli di uno spettacolo.
            \item \textbf{For} ogni eventi vincolato allo spettacolo
            \begin{enumerate}
                \item Il Sistema mostra informazioni rilevanti dell'evento (es. data, ora, sala).
            \end{enumerate}
        \end{enumerate} \\
        \hline
        \cellcolor{gray!20}
        \textbf{Postcondizioni:} \\
        \cellcolor{gray!20}
        \begin{minipage}{\linewidth}
            \begin{enumerate}
                \item E' visualizzata a schermo la lista degli eventi associati allo spettacolo.
            \end{enumerate}
        \end{minipage} \\
        \hline
        \textbf{Sequenza degli eventi alternativa:} \\
        Nessuna. \\
        \hline
    \end{tabular}
\end{table}

% UC: Lista Sezioni
\begin{table}[H]
    \begin{tabular}{|p{\linewidth}|}
        \hline
        \cellcolor{gray!100}
        \color{white}
        \centerline{\textbf{Caso d'uso:} Lista Sezioni} \\ 
        \hline
        \textbf{ID:} 9 \\
        \hline
        \cellcolor{gray!20}
        \textbf{Breve descrizione:} \\
        \cellcolor{gray!20}
        Messa a schermo delle sezioni. \\
        \hline
        \textbf{Attori primari:} \\
        \begin{minipage}{\linewidth}
            Cliente \\
            Biglietteria \\
            Amministratore 
        \end{minipage}
        \vspace{0pt} \\  % Utilizzo \vspace{} per regolare lo spazio tra le sezioni delle tabelle.
        \hline
        \textbf{Attori secondari:} \\
        Nessuno. \\
        \hline
        \cellcolor{gray!20}
        \textbf{Precondizioni:} \\
        \cellcolor{gray!20}
        \begin{minipage}{\linewidth}
            \begin{enumerate}
            	\item L'utente si trova nella categoria "Sezioni \& Posti"
                \item L'utente sta visualizzando l'elenco delle sezioni.
            \end{enumerate}
        \end{minipage} 
        \vspace{0pt} \\
        \hline
        \textbf{Sequenza degli eventi principale:}
        \begin{enumerate}
            \item Il caso d'uso inizia quando l'utente seleziona \textit{Sezioni}.
            \item \textbf{include}(Visualizza Sezione)
            \item \textbf{If} esistono sezioni
            \begin{enumerate}
                \item Il Sistema mostra all'utente tutte le Sezioni presenti una dopo l'altra.
            \end{enumerate}
            \item \textbf{Else}
            \begin{enumerate}
                \item Il Sistema comunica che la lista delle Sezioni è vuota.
                \item Il Sistema mostra un elenco vuoto.
            \end{enumerate}				
        \end{enumerate} \\
        \hline
        \cellcolor{gray!20}
        \textbf{Postcondizioni:} \\
        \cellcolor{gray!20}
        \begin{minipage}{\linewidth}
            \begin{enumerate}
                \item Il Sistema ha elencato le Sezioni presenti.
            \end{enumerate}
        \end{minipage} \\
        \hline
        \textbf{Sequenza degli eventi alternativa:} \\
        Nessuna. \\
        \hline
    \end{tabular}
\end{table}

% UC: Visualizza Sezione
\begin{table}[H]
    \begin{tabular}{|p{\linewidth}|}
        \hline
        \cellcolor{gray!100}
        \color{white}
        \centerline{\textbf{Caso d'uso:} Visualizza Sezione} \\
        \hline
        \textbf{ID:} 10 \\
        \hline
        \cellcolor{gray!20}
        \textbf{Breve descrizione:} \\
        \cellcolor{gray!20}
        Visualizzazione con dettagli di una sezione. \\
        \hline
        \textbf{Attori primari:} \\
        \begin{minipage}{\linewidth}
            Amministratore \\
            Biglietteria \\
            Cliente
        \end{minipage}
        \vspace{0pt} \\ 
        \hline
        \textbf{Attori secondari:} \\
        Nessuno. \\
        \hline
        \cellcolor{gray!20}
        \textbf{Precondizioni:} \\
        \cellcolor{gray!20}
        \begin{minipage}{\linewidth}
            \begin{enumerate}
                \item L'utente sta ha richiesto la lista delle sezioni.
            \end{enumerate}
        \end{minipage} \\
        \hline
        \textbf{Sequenza degli eventi principale:}
        \begin{enumerate}
            \item Il caso d'uso inizia quando l'utente visualizza una Sezione.
            \item \textbf{include}(Lista Posti)
            \item Il Sistema mostra all'utente i dati della Sezione.
            \item Il Sistema mostra il pulsante \textit{Lista Posti}.
            \item \textbf{If} l'utente è autenticato come Amministratore
            \begin{enumerate}
                \item Il Sistema mostra, accanto al primo pulsante, altre due opzioni: \textit{Modifica Sezione} e \textit{Elimina Sezione}.
                \item[] \textbf{punto di estensione}: modificaSezione, eliminaSezione
            \end{enumerate}
        \end{enumerate} \\
        \hline
        \cellcolor{gray!20}
        \textbf{Postcondizioni:} \\
        \cellcolor{gray!20}
        \begin{minipage}{\linewidth}
            \begin{enumerate}
                \item Il Sistema ha mostrato i dati della Sezione.
                \item Il Sistema dà l'opzione di visualizzare i posti relativi alla Sezione.
                \item Se l'utente è Amministratore, il Sistema permette di modificare o eliminare la sezione.
            \end{enumerate}
        \end{minipage} \\
        \hline
        \textbf{Sequenza degli eventi alternativa:} \\
        Nessuna. \\
        \hline
    \end{tabular}
\end{table}

% UC: Lista Posti
\begin{table}[H]
    \begin{tabular}{|p{\linewidth}|}
        \hline
        \cellcolor{gray!100}
        \color{white}
        \centerline{\textbf{Caso d'uso:} Lista Posti} \\
        \hline
        \textbf{ID:} 11 \\
        \hline
        \cellcolor{gray!20}
        \textbf{Breve descrizione:} \\
        \cellcolor{gray!20}
        Messa a schermo dei posti. \\
        \hline
        \textbf{Attori primari:} \\
        \begin{minipage}{\linewidth}
            Amministratore \\
            Biglietteria \\ 
            Cliente
        \end{minipage}
        \vspace{0pt} \\ 
        \hline
        \textbf{Attori secondari:} \\
        Nessuno. \\
        \hline
        \cellcolor{gray!20}
        \textbf{Precondizioni:} \\
        \cellcolor{gray!20}
        \begin{minipage}{\linewidth}
            \begin{enumerate}
                \item È stato caricato in memoria un database valido. 
            \end{enumerate}
        \end{minipage} 
        \vspace{-10pt} \\
        \hline
        \textbf{Sequenza degli eventi principale:}
        \begin{enumerate}
            \item Il caso d'uso inizia quando l'utente seleziona \textit{Posti} all'interno di una Sezione.
            \item \textbf{If} la Sezione contiene Posti
            \begin{enumerate}
                \item Il Sistema mostra all'utente tutti i Posti presenti della Sezione sotto forma di elenco.
            \end{enumerate}
            \item \textbf{Else}
            \begin{enumerate}
                \item Il Sistema comunica che la lista dei Posti della Sezione è vuota.
                \item Il Sistema mostra un elenco vuoto.
            \end{enumerate}				
        \end{enumerate} \\
        \hline
        \cellcolor{gray!20}
        \textbf{Postcondizioni:} \\
        \cellcolor{gray!20}
        \begin{minipage}{\linewidth}
            \begin{enumerate}
                \item Il Sistema ha elencato i Posti della Sezione selezionata.
            \end{enumerate}
        \end{minipage} \\
        \hline
        \textbf{Sequenza degli eventi alternativa:} \\
        Nessuna. \\
        \hline
    \end{tabular}
\end{table}

% UC: Crea Prenotazione
\begin{table}[H]
    \begin{tabular}{|p{\linewidth}|}
        \hline
        \cellcolor{gray!100}
        \color{white}
        \centerline{\textbf{Caso d'uso:} Crea Prenotazione} \\
        \hline
        \textbf{ID:} 12 \\
        \hline
        \cellcolor{gray!20}
        \textbf{Breve descrizione:} \\
        \cellcolor{gray!20}
        L'utente prenota posti di un evento svolto nel teatro. \\
        \hline
        \textbf{Attori primari:} \\
        \begin{minipage}{\linewidth}
            Cliente \\
            Biglietteria \\
            Amministratore
        \end{minipage}
        \vspace{0pt} \\
        \hline
        \textbf{Attori secondari:} \\                        
        Nessuno. \\
        \hline
        \cellcolor{gray!20}
        \textbf{Precondizioni:} \\
        \cellcolor{gray!20}
        \begin{minipage}{\linewidth}
            \begin{enumerate}
                \item L'utente visualizza l'elenco degli eventi.
            \end{enumerate}
        \end{minipage} \\
        \hline
        \textbf{Sequenza degli eventi principale:}
        \begin{enumerate}
            \item Il caso d'uso inizia quando l'utente preme il pulsante di creazione.
            \item Il Sistema, tramite un'interfaccia, permette di inserire i dati relativi alla prenotazione.
            \item L'utente inserisce i dati relativi alla prenotazione.
            \item L'utente salva i dati.
            \item Il Sistema verifica i dati e li salva in memoria.
        \end{enumerate} \\
        \hline
        \cellcolor{gray!20}
        \textbf{Postcondizioni:} \\
        \cellcolor{gray!20}
        \begin{minipage}{\linewidth}
            \begin{enumerate}
                \item L'utente crea una prenotazione.
            \end{enumerate}
        \end{minipage}
        \vspace{-10pt} \\
        \hline
        \textbf{Sequenza degli eventi alternativa:} \\
        Nessuna. \\
        \hline
    \end{tabular}
\end{table}

% UC: Visualizza Prenotazione
\begin{table}[H]
    \begin{tabular}{|p{\linewidth}|}
        \hline
        \cellcolor{gray!100}
        \color{white}
        \centerline{\textbf{Caso d'uso:} Visualizza Prenotazione} \\
        \hline
        \textbf{ID:} 13 \\
        \hline
        \cellcolor{gray!20}
        \textbf{Breve descrizione:} \\
        \cellcolor{gray!20}
        L'utente ottiene l'informazione anagrafica riguardo una singola prenotazione. \\
        \hline
        \textbf{Attori primari:} \\
        \begin{minipage}{\linewidth}
            Cliente \\
            Biglietteria \\
            Amministratore
        \end{minipage}
        \vspace{0pt} \\
        \hline
        \textbf{Attori secondari:} \\
        Nessuno. \\
        \hline
        \cellcolor{gray!20}
        \textbf{Precondizioni:} \\
        \cellcolor{gray!20}
        \begin{minipage}{\linewidth}
            \begin{enumerate}
                \item L'utente visualizza la lista delle prenotazioni.
                \item L'elenco delle prenotazioni, sia esso filtrato o meno, non è vuoto.
            \end{enumerate}
        \end{minipage}
        \vspace{-5pt} \\
        \hline
        \textbf{Sequenza degli eventi principale:}
        \begin{enumerate}
            \item Il caso d'uso inizia quando l'utente richiede le informazioni di una prenotazione.
            \item Il Sistema mostra le informazioni relative alla prenotazione.
            \item Il Sistema mostra tre pulsanti per modificare, eliminare o generare la ricevuta della prenotazione.
            \item \textbf{If} l'utente è Amministratore o Biglietteria
            \begin{enumerate}
                \item Il Sistema mostra un pulsante per segnare la prenotazione come \emph{pagata}.
            \end{enumerate}
            \item[] \textbf{punto di estensione:} modificaPrenotazione
        \end{enumerate} \\
        \hline
        \cellcolor{gray!20}
        \textbf{Postcondizioni:} \\
        \cellcolor{gray!20}
        \begin{minipage}{\linewidth}
            \begin{enumerate}
                \item L'utente visualizza le informazioni di una prenotazione.
                \item Il Sistema permette all'utente di modificare, eliminare o generare una ricevuta della prenotazione.
                \item Se l'utente è Amministratore o Biglietteria, il Sistema permette di segnarla come \emph{pagata}.
            \end{enumerate}
        \end{minipage}
        \vspace{0pt} \\
        \hline
        \textbf{Sequenza degli eventi alternativa:} \\
        Nessuna. \\
        \hline
    \end{tabular}
\end{table}

% UC EXT: Modifica Prenotazione
\begin{table}[H]
    \begin{tabular}{|p{\linewidth}|}
        \hline
        \cellcolor{gray!100}
        \color{white}
        \centerline{\textbf{Caso d'uso di estensione:} Modifica Prenotazione} \\
        \hline
        \textbf{ID:} 14 \\
        \hline
        \cellcolor{gray!20}
        \textbf{Breve descrizione:} \\
        \cellcolor{gray!20}
        \textbf{Segmento 1:} L'utente modifica una prenotazione esistente all'interno dell'elenco delle prenotazioni. \\
        \hline
        \textbf{Attori principali:} \\
        \begin{minipage}{\linewidth}
            Cliente \\
            Biglietteria \\
            Amministatore
        \end{minipage}
        \vspace{0pt} \\
        \hline
        \textbf{Attori secondari:} \\                        
        Nessuno. \\
        \hline
        \cellcolor{gray!20}
        \textbf{Precondizioni del segmento 1:} \\
        \cellcolor{gray!20}
        \begin{minipage}{\linewidth}
            \begin{enumerate}
                \item L'utente visualizza l'elenco delle prenotazioni.
                \item L'utente ha premuto il pulsante di modifica.
            \end{enumerate}
        \end{minipage}
        \vspace{-5pt} \\
        \hline
        \textbf{Sequenza degli eventi del segmento 1:}
        \begin{enumerate}
            \item Il Sistema mostra la interfaccia di creazione di prenotazione, ma con i dati saltavi della prenotazione già creata selezionati.
            \item L'utente modifica i dati relativi alla prenotazione.
            \item L'utente salva i dati.
            \item Il Sistema verifica i dati e li salva in memoria.
        \end{enumerate} \\
        \hline
        \cellcolor{gray!20}
        \textbf{Postcondizioni del segmento 1:} \\
        \cellcolor{gray!20}
        \begin{minipage}{\linewidth}
            \begin{enumerate}
                \item L'utente effettua la modifica di una prenotazione.
            \end{enumerate}
        \end{minipage} \\
        \hline
    \end{tabular}
\end{table}

% UC EXT: Elimina Prenotazione
\begin{table}[H]
    \begin{tabular}{|p{\linewidth}|}
        \hline
        \cellcolor{gray!100}
        \color{white}
        \centerline{\textbf{Caso d'uso di estensione:} Elimina Prenotazione} \\
        \hline
        \textbf{ID:} 15 \\
        \hline
        \cellcolor{gray!20}
        \textbf{Breve descrizione:} \\
        \cellcolor{gray!20}
        \textbf{Segmento 1:} L'utente elimina una prenotazione esistente all'interno dell'elenco delle prenotazioni. \\
        \hline
        \textbf{Attori principali:} \\
        \begin{minipage}{\linewidth}
            Cliente \\
            Biglietteria \\
            Amministratore
        \end{minipage}
        \vspace{0pt} \\
        \hline
        \textbf{Attori secondari:} \\                        
        Nessuno. \\
        \hline
        \cellcolor{gray!20}
        \textbf{Precondizioni del segmento 1:} \\
        \cellcolor{gray!20}
        \begin{minipage}{\linewidth}
            \begin{enumerate}
                \item L'utente visualizza l'elenco delle prenotazioni.
                \item L'utente ha premuto il pulsante di eliminazione.
            \end{enumerate}
        \end{minipage}
        \vspace{-5pt} \\
        \hline
        \textbf{Sequenza degli eventi del segmento 1:}
        \begin{enumerate}
            \item Il Sistema chiede conferma per eliminare la prenotazione.
            \item L'utente conferma l'eliminazione della prenotazione.
            \item Il Sistema effettua l'eliminazione.
        \end{enumerate} \\
        \hline
        \cellcolor{gray!20}
        \textbf{Postcondizioni del segmento 1:} \\
        \cellcolor{gray!20}
        \begin{minipage}{\linewidth}
            \begin{enumerate}
                \item L'utente effettua l'eliminazione di una prenotazione.
            \end{enumerate}
        \end{minipage} \\
        \hline
    \end{tabular}
\end{table}

% UC: Genera Ricevuta
\begin{table}[H]
    \begin{tabular}{|p{\linewidth}|}
        \hline
        \cellcolor{gray!100}
        \color{white}
        \centerline{\textbf{Caso d'uso:} Genera Ricevuta} \\
        \hline
        \textbf{ID:} 16 \\
        \hline
        \cellcolor{gray!20}
        \textbf{Breve descrizione:} \\
        \cellcolor{gray!20}
        L'utente genera la ricevuta di una prenotazione e la stampa. \\
        \hline
        \textbf{Attori primari:} \\
        \begin{minipage}{\linewidth}
            Cliente \\
            Biglietteria \\
            Amministratore
        \end{minipage}
        \vspace{0pt} \\
        \hline
        \textbf{Attori secondari:} \\
        Nessuno. \\
        \hline
        \cellcolor{gray!20}
        \textbf{Precondizioni:} \\
        \cellcolor{gray!20}
        \begin{minipage}{\linewidth}
            \begin{enumerate}
                \item L'utente visualizza una prenotazione. %/""
            \end{enumerate}
        \end{minipage} \\
        \hline
        \textbf{Sequenza degli eventi principale:}
        \begin{enumerate}
            \item Il caso d'uso inizia quando l'utente preme \emph{Genera ricevuta}. %/""
            \item Il Sistema genera un documento con tutta l'informazione relativa all'acquisto dei biglietti prenotati.
            \item Il Sistema stampa il documento.
        \end{enumerate} \\
        \hline
        \cellcolor{gray!20}
        \textbf{Postcondizioni:} \\
        \cellcolor{gray!20}
        \begin{minipage}{\linewidth}
            \begin{enumerate}
                \item La ricevuta è stata stampata.
            \end{enumerate}
        \end{minipage} \\
        \hline
        \textbf{Sequenza degli eventi alternativa:} \\
        Nessuna. \\
        \hline
    \end{tabular}
\end{table}
